% 第 3 章 全球存储产业链图谱
% 含 mermaid 产业链架构图 (tikz 等效实现) + 环节价值/毛利/对标表

\section{产业链七层架构与价值分布}

全球存储产业链从上游到下游可划分为七层：上游设备与材料、存储设计、晶圆制造、封装 (先进+传统)、模组与主控、整机集成、下游应用 (AI/云/终端)。A 股映射的核心机会集中在设备材料（国产化率 10-15\%，翻倍空间）、接口芯片设计（澜起 RCD ≈70\% + CXL 首发卡位）与先进封装（长电\slash 通富 弹性）三层。

\needspace{52\baselineskip}
\begin{exhibitbox}[图 3-1 \quad 全球存储产业链全景图谱 / Global Memory Supply-Chain Architecture]
\centering
\resizebox{0.88\exhibitwidth}{!}{%
\begin{tikzpicture}[
  >=Stealth,
  node distance=0.55cm,
  layer/.style={rectangle, rounded corners=3pt, draw=#1, fill=#1!8, text width=13cm, align=center, inner sep=3pt, font=\bfseries\small, minimum height=0.55cm},
  node/.style={rectangle, rounded corners=2pt, draw=steelgrey, fill=white, text width=2.35cm, align=center, minimum height=0.95cm, font=\scriptsize, inner sep=2pt},
  ashar/.style={rectangle, rounded corners=2pt, draw=riskamber, fill=white, text width=2.35cm, align=center, minimum height=0.95cm, font=\scriptsize\bfseries, text=deepnavy, inner sep=2pt, line width=0.7pt},
  gap/.style={text width=0.3cm}
]
% Layer 1
\node[layer=accentblue] (L1) {\faIndustry\enspace L1 · 上游 · 设备与材料 \quad[\;设备 10--12\% \;|\; 材料 7--9\%\;]};
\node[node, below=of L1.west, anchor=north west, xshift=0.25cm] (U1) {EUV/光刻\\ASML · Nikon\\\textcolor{mutedgrey}{国产化 <1\%}};
\node[node, right=of U1] (U2) {刻蚀/薄膜/ALD\\Lam · AMAT · TEL\\\textcolor{mutedgrey}{国产化 10--15\%}};
\node[node, right=of U2] (U3) {测试设备\\Advantest · Teradyne\\\textcolor{mutedgrey}{国产化 5--8\%}};
\node[node, right=of U3] (U4) {硅片/光刻胶/特气\\ShinEtsu · JSR · Entegris\\\textcolor{mutedgrey}{国产化 8--12\%}};
\node[ashar, fill=white, draw=riskamber, text=deepnavy, right=of U4] (A1) {\faStar\enspace A 股 · 设备\\北方华创 · 中微\\拓荆 · 沪硅 · 南大};

% Layer 2
\node[layer=nvgreen, below=1.1cm of L1] (L2) {\faCogs\enspace L2 · 存储设计 \quad[\;8--10\%\;]};
\node[node, below=of L2.west, anchor=north west, xshift=0.25cm] (D1) {DRAM 设计\\Samsung · SK · Micron\\\textcolor{mutedgrey}{国产 <1\%}};
\node[node, right=of D1] (D2) {NAND 设计\\Samsung · Kioxia · YMTC\\\textcolor{mutedgrey}{国产 1--3\%}};
\node[node, right=of D2] (D3) {接口/主控/RCD\\Rambus · Marvell\\\textcolor{mutedgrey}{澜起领先}};
\node[ashar, fill=white, draw=riskamber, text=deepnavy, right=of D3] (A2) {\faStar\enspace A 股 · 设计\\澜起科技 (RCD 70\%)\\兆易 · 芯原 · 北京君正};

% Layer 3
\node[layer=riskamber, below=1.1cm of L2] (L3) {\faMicrochip\enspace L3 · 晶圆制造 \quad[\;38--42\%\;] \textcolor{riskred}{—— 价值最大环节}};
\node[node, below=of L3.west, anchor=north west, xshift=0.25cm] (M1) {DRAM FAB\\SK · Samsung · Micron\\CXMT (长鑫, 未上市)};
\node[node, right=of M1] (M2) {NAND FAB\\Samsung · Kioxia · Micron\\YMTC (长存, 未上市)};
\node[ashar, fill=white, draw=riskamber, text=deepnavy, right=of M2] (A3) {\faStar\enspace A 股间接\\\small A 股无直接原厂\\北方华创等为供应商};

% Layer 4
\node[layer=riskred, below=1.1cm of L3] (L4) {\faLayerGroup\enspace L4 · 封装 \quad[\;先进 12--15\% \;|\; 传统 5--7\%\;]};
\node[node, below=of L4.west, anchor=north west, xshift=0.25cm] (P1) {CoWoS/TSV 先进\\TSMC (约80\%) · Samsung\\ASE · Amkor};
\node[node, right=of P1] (P2) {传统 FBGA/LGA\\ASE · Amkor · SPIL\\\textcolor{mutedgrey}{国产 约30\%}};
\node[ashar, fill=white, draw=riskamber, text=deepnavy, right=of P2] (A4) {\faStar\enspace A 股 · 封测\\长电 (XDFOI)\\通富 (AMD 绑定) · 华天};

% Layer 5
\node[layer=deepnavy, below=1.1cm of L4] (L5) {\faMemory\enspace L5 · 模组/主控 \quad[\;8--10\%\;]};
\node[node, below=of L5.west, anchor=north west, xshift=0.25cm] (MO1) {企业级模组/SSD\\Smart Global · Dell\\Samsung OEM};
\node[node, right=of MO1] (MO2) {消费级模组\\Apacer · G.Skill\\Kingston (深科技代工)};
\node[ashar, fill=white, draw=riskamber, text=deepnavy, right=of MO2] (A5) {\faStar\enspace A 股 · 模组\\江波龙 (Foresee\/Lexar)\\深科技 (沛顿) · 佰维存储};

% Layer 6
\node[layer=navy, below=1.1cm of L5] (L6) {\faServer\enspace L6 · 整机集成 \quad[\;-\;]};
\node[node, below=of L6.west, anchor=north west, xshift=0.25cm] (S1) {AI 服务器\\NVIDIA DGX · Dell\\联想 · 浪潮 · 华为};
\node[node, right=of S1] (S2) {PC/手机/汽车\\Apple · Dell · HP\\比亚迪 · 小米};
\node[ashar, fill=white, draw=riskamber, text=deepnavy, right=of S2] (A6) {\faStar\enspace A 股 · 系统\\(本报告不覆盖)\\浪潮 · 联想 · 华为};

% Layer 7
\node[layer=steelgrey, below=1.1cm of L6] (L7) {\faGlobe\enspace L7 · 下游应用 (Capex 来源)};
\node[node, fill=steelgrey!10, below=of L7.west, anchor=north west, xshift=0.25cm, text width=4.1cm] (C1) {AI 训练集群\\MSFT · Meta · Google · AWS\\字节 · 腾讯 · 阿里 (\$2500 亿+)};
\node[node, fill=steelgrey!10, right=of C1, text width=4.1cm] (C2) {云厂商数据中心\\AWS · Azure · GCP\\阿里云 · 腾讯云 · 火山};
\node[node, fill=steelgrey!10, right=of C2, text width=4.1cm] (C3) {消费电子 + 汽车\\\(> \$10000 亿/年\) TAM\\PC 复苏 + AI PC + 新能源};

% Main arrows
\foreach \i in {1,...,6} {
  \pgfmathsetmacro\j{int(\i+1)};
  \draw[->, thick, navy] (L\i.south) -- node[midway, right=1pt, font=\tiny\color{steelgrey}] {价值传导} (L\j.north);
}

% Key stat
\node[draw=navy, fill=white, rounded corners=3pt, font=\scriptsize\bfseries, text=navy, right=0.2cm of U2, yshift=-0.55cm] {\faExclamationTriangle A 股无原厂 · 配置聚焦上游/接口/封测};
\end{tikzpicture}
}
\end{exhibitbox}
\sourcenote{产业链架构设计参考 Industry Landscape §1.3 mermaid 原版；价值占比来自中金 2026 年 Q1 半导体深度报告 + AStock 交叉验证。}

\vspace{0.4em}
\noindent\textbf{投资含义}：产业链七层价值分布呈现显著的「哑铃型」——晶圆制造 (38-42\%) 和设备材料 (17-21\%) 合计占据近六成价值。A 股恰恰在晶圆制造层缺失 (YMTC\slash CXMT 未上市)，但在设备材料和先进封装两个高价值增量层拥有全球竞争力的标的。这一结构决定了 A 股存储链投资必须走「上游设备材料（国产化率 10-15\%，18 个月翻倍空间）+ 接口芯片（澜起 RCD 全球 ≈70\% 份额）」的路径，而非追逐不存在的「国产 HBM 原厂」概念。

\section{环节价值量、毛利率、A 股对标与国产化率四维总表}

\begin{exhibitbox}[表 3-1 \quad 产业链各环节价值量 / 毛利率 / A 股对标 / 国产化率 / 技术代差 五维总表]
\centering
\footnotesize
\begin{tabularx}{\exhibitboxwidth}{>{\raggedright}p{1.5cm} >{\raggedleft}p{0.85cm} >{\raggedleft}p{0.85cm} >{\raggedright\arraybackslash\hspace{0pt}}p{2.0cm} >{\raggedleft}p{0.85cm} >{\raggedright\arraybackslash\hspace{0pt}}p{1.4cm} >{\raggedright\arraybackslash\hspace{0pt}}X}
\toprule
\textbf{环节} & \textbf{价值(\%)} & \textbf{毛利(\%)} & \textbf{全球龙头} & \textbf{国产化(\%)} & \textbf{A 股核心标的} & \textbf{A 股卡位与代差评估} \\
\midrule
\rowcolor{accentblue!5}
\textbf{半导体设备} & 10--12 & 45--60 & Lam · AMAT · ASML & {\faExclamationTriangle} 10--15 & 北华创\newline 中微公司\newline 拓荆科技 & 刻蚀/薄膜/ALD 在成熟制程全面突破，<16nm 验证中；\textbf{二阶 Beta 首选层}，国产化率仍有翻倍空间。代差 1--2 代。 \\
\midrule
\rowcolor{accentblue!5}
\textbf{半导体材料} & 7--9 & 35--50 & ShinEtsu · JSR & {\faExclamationTriangle} 8--12 & 沪硅产业\newline 南大光电 & 12 寸硅片 28nm 批量供应；ArF 光刻胶 2026 目标突破；KrF/特气已破。代差 1--3 代。 \\
\midrule
\rowcolor{nvgreen!5}
\textbf{存储设计·接口} & 8--10 & 55--70 & Rambus · Marvell & 约3--5 & 澜起科技 ★ & \textbf{全球领先无代差}——RCD 全球 70\% 份额，CXL 首发；A 股唯一「一阶 Alpha」。兆易 NOR 前三，DRAM 追赶。代差 0。 \\
\midrule
\textbf{存储设计·原厂} & 8--10 & 40--60 & SK · Samsung · Micron & <1 & 无直接标的 & YMTC\slash CXMT 未上市；兆易自研 DRAM 试产；原厂层缺失是最大结构缺口。代差 3--4 代 (HBM)。 \\
\midrule
\textbf{晶圆制造} & 38--42 & 30--50 & Samsung · SK · YMTC & {\faExclamationTriangle} 8--15 & 无直接标的 & 长存\slash 长鑫未上市，A 股仅通过设备商间接映射；DRAM 代差 1.5--2 代，NAND 代差 1 代。 \\
\midrule
\rowcolor{riskred!5}
\textbf{先进封装 (CoWoS)} & 12--15 & 25--35 & TSMC (约80\%) & {\faExclamationTriangle} 5--8 & 长电科技\newline 通富微电 & TSMC 垄断 HBM 封装产能；长电 XDFOI/通富 CoWoS 2027 前难撼份额；\textbf{弹性标的}，AI 实际贡献 <3\%。代差 1--2 代。 \\
\midrule
\textbf{传统封测} & 5--7 & 15--22 & ASE · Amkor & 约25--30 & 长电\slash 通富/华天 & 传统封装国产替代较充分，毛利率低。代差 <1 代。 \\
\midrule
\rowcolor{deepnavy!5}
\textbf{存储模组/SSD} & 8--10 & 10--20 & Smart Global · ATP & {\faExclamationTriangle} 10--20 & 江波龙\newline 深科技 & 深科技 Kingston 代工确定性高；江波龙 Foresee 企业级 SSD 导入头部云厂。代差 0--1 代。 \\
\bottomrule
\end{tabularx}
\end{exhibitbox}
\sourcenote{龙头毛利率来自各公司 FY2025 年报 (海外原厂/设备)；国产化率综合中国国际招标网公开数据 + {\faExclamationTriangle} 行业推断；价值占比参考中金半导体深度 [L3]。}

\vspace{0.4em}
\noindent\textbf{投资含义}：五维对比揭示三个关键投资洞察。\textbf{(1) 最佳配置层：设备材料 + 接口芯片。}设备层国产化率 10-15\% (翻倍空间) + 全球 capex 上行周期的「双重 Beta」叠加上升；接口层 (澜起) 是 A 股唯一全球领先无代差的环节。\textbf{(2) 需规避的结构性缺失层：原厂设计 + 晶圆制造。}A 股真正的原厂缺失意味着国产 HBM 主题标的本质上是「映射的映射」，估值弹性大但基本面脆弱。\textbf{(3) 先进封装是「远期弹性层」而非「当前基本面层」。}长电\slash 通富当前 AI 收入实际贡献 <3\%，需 2027 年 H2+ 方能体现 CoWoS 产能贡献，当前估值已部分反映远期预期，需警惕节奏误判。
